\documentstyle[% 


righttag% 


,11pt% 


,amssymb% 


,russian%               remove in Eng. 


,izvru%                 Set izven% in Eng. 


% ,izvru-ks             for brief communications 


]{amsart} 


 


%      Math definitions 


 


\newcommand{\la}{\langle} 


\newcommand{\ra}{\rangle} 


\newcommand{\mes}{\operatorname{mes}} 


\newcommand{\supess}{\operatornamewithlimits{sup\,ess}} 


\newcommand{\tts}[1]{} 


 


\theoremstyle{plain} 


\theorembodyfont{\it} 


\newtheorem{thmnonum}{\hskip\parindent Теорема} 


\renewcommand{\thethmnonum}{} 


 


\theorembodyfont{\rm} 


\newtheorem{notenonum}{\hskip\parindent Замечание} 


\newtheorem{defnnonum}{\hskip\parindent Определение} 


\renewcommand{\thenotenonum}{} 


\renewcommand{\thedefnnonum}{} 


\numberwithin{equation}{section} 


 


%\hoffset=-15mm%                  For Eng. 


  


\setcounter{page}{37} 


\god{1998} 


\nomer{10~(437)} %                        Remove in Eng. 


%\nomer{Vol.\,42, No.\,10}%               Set in Eng. 


%\str{\pageref{begin}--\pageref{end}}%  Set in Eng. 


\udk{517.988.8} 


\author{\it С.Е.\,Железовский} 


% NB:   ^^ remove in Eng. 


\title{О существовании и единственности решения\\ 


и о скорости сходимости метода Бубнова--Гал\"еркина\\ 


для одной квазилинейной эволюционной задачи\\ 


в гильбертовом пространстве} 


 


\date{} 


% \pagestyle{myheadings}%         Set in Eng. 


 


\pagestyle{plain} %              Remove in Eng. 


\begin{document} 


 


% \thispagestyle{plain}%          Set in Eng. 


 


\maketitle 


\label{begin} 


 


\section{Введение} 


 


В данной статье рассматривается задача Коши для квазилинейного 


эволюционного уравнения второго порядка:


\begin{gather} 


Q[u,t]\equiv u''(t)+K(t)u'(t)+A(t)u(t)+B(t,u(t))=0,\quad 


t\in[0,T], \\ 


u(0)=u^{(0)}_0,\quad u'(0)=u^{(0)}_1. 


\end{gather} 


Здесь $u$ --- искомая функция, отображающая отрезок $[0,T]$ числовой оси 


$\Bbb R$ в 


вещественное сепарабельное гильбертово пространство $H$; $K(t)$, $A(t)$ 


($t\in[0,T]$) 


--- линейные операторы, действующие в $H$, причем все операторы $A(t)$ 


самосопряженные и положительно-определенные; $B$ --- нелинейный 


оператор; $u^{(0)}_0$, $u^{(0)}_1$ --- заданные элементы пространства 


$H$. Условия на 


операторные коэффициенты и начальные данные задачи (1.1), (1.2) 


подробно сформулированы ниже в п.\,2. В класс задач, для которых эти 


условия выполнены, входят, в частности, некоторые начально-краевые 


задачи математической физики, получаемые на основе вариационного 


принципа Гамильтона--Остроградского. 


 


Основной целью статьи является получение для задачи (1.1), (1.2)  


априорной оценки погрешности метода Бубнова--Гал\"еркина (МБГ) при 


специальном выборе базиса (в качестве базисных элементов МБГ 


используются собственные элементы вспомогательного оператора, сходного 


(\cite{1}, с.\,21) и образующего острый угол (\cite{2}) с каждым 


из операторов $A(t)$\,). Оценки такого типа ранее были получены для 


эволюционных задач, содержащих вторую производную искомой функции по 


времени, в работах \cite{3}--\cite{5}, для стационарных задач и 


эволюционных задач первого порядка по времени -- в ряде работ, из 


которых отметим \cite{6}--\cite{9}. 


 


Наряду с оценкой погрешности МБГ для задачи (1.1), (1.2) 


устанавливается результат об однозначной разрешимости, необходимый для 


полноты и корректности изложения. В имеющейся литературе такие 


результаты для рассматриваемой задачи или для задач более общего вида 


отсутствуют. Для нелинейных задач в абстрактном гильбертовом 


пространстве, подобных рассматриваемой, результат о существовании 


обобщенного решения получен в \cite{10} и результат об устойчивости, из 


которого следует единственность решения, --- в \cite{11}. 


 


Отметим, что метод доказательства существования решения, принятый в 


данной работе, отличается от известного метода компактности, обычно 


используемого совместно с МБГ при исследовании задач, подобных (1.1), 


(1.2) (см., напр., \cite{10}, а также \cite{12}, гл.\,I). Именно, в 


нашей работе устанавливается априорная оценка разности между любыми 


двумя приближенными решениями рассматриваемой задачи, построенными по 


МБГ, доставляющая результат о сильной сходимости последовательности 


приближенных решений и соответственно о существовании точного решения. 


Таким образом, доказательство существования решения совмещается с 


выводом оценки погрешности МБГ, которая в данном случае получается 


предельным переходом в оценке разности между двумя приближенными 


решениями. Напомним, что при использовании обычного метода компактности 


непосредственно устанавливается только слабая сходимость 


последовательности приближенных решений или некоторой ее 


подпоследовательности. 


 


Пусть $X$ и $Y$ --- вещественные банаховы пространства с нормами 


$|\cdot|_X$ и $|\cdot|_Y$ 


соответственно. В дальнейшем используются следующие обозначения: 


$\cal L(X,Y)$ --- 


пространство линейных ограниченных операторов $L:X\to Y$ с нормой 


$\|L\mid\cal L(X,Y)\|= 


\sup\limits_{u\in X\setminus\{0\}}|Lu|_Y|u|^{-1}_X$; 


$\cal L(X)=\cal L(X,X)$; $\cal B(X)$ 


--- пространство симметричных билинейных форм $b:X^2\to\Bbb R$, 


удовлетворяющих 


условию 


$\|b\mid\cal B(X)\|\equiv 


\sup\limits_{u,v\in X\setminus\{0\}}b(u,v)|u|^{-1}_X|v|^{-1}_X<\infty$; 


$L^p(0,T;X)$ ($1\le p\le\infty$, $T>0$) 


--- пространство сильно измеримых 


функций $u:[0,T]\to X$, удовлетворяющих условию 


$$ 


\|u\mid L^p(0,T;X)\|\equiv\biggl(\,\int^T_0|u(t)|^p_X dt 


\biggr)^{1/p}<\infty 


$$ 


(при $p=\infty$ вместо этого требуется 


$\|u\mid L^\infty(0,T;X)\|\equiv 


\supess\limits_{0\le t\le T}|u(t)|_X<\infty$); 


$L(0,T;X)=L^1(0,T;X)$; $L(0,T)=L(0,T;\Bbb R)$; $C([0,T];X)$ 


--- пространство функций 


$u:[0,T]\to X$, непрерывных на $[0,T]$, с нормой 


$\|u\mid C([0,T];X)\|=\max\limits_{0\le t\le T}|u(t)|_X$; 


$C([0,T])=C([0,T];\Bbb R)$; $C^k([0,T];X)$ ($k\in\Bbb N$) --- 


пространство 


функций, отображающих $[0,T]$ в $X$ и $k$ раз непрерывно 


дифференцируемых на 


$[0,T]$. Производные, не существующие в обычном смысле, понимаются в 


данной работе согласно следующему {\it определению}:


функция $v:[0,T]\to X$ называется 


обобщенной производной функции $u:[0,T]\to X$ (т.\,е.~$v=u'$), если она 


принадлежит $L(0,T;X)$ 


и $u(t)=u(0)+\int\limits^t_0{}v(\tau)d\tau$ для всех $t\in[0,T]$, 


$u\in C(0,T;H)$, где интеграл понимается в смысле Бохнера. 


Запись $u'\in L^p(0,T;X)$ всюду ниже означает, что обобщенная 


производная функции $u:[0,T]\to X$ 


существует и принадлежит $L^p(0,T;X)$ (в частности, запись 


$u'\in L(0,T;X)$ означает 


только то, 


что обобщенная производная $u'$ существует). Через $\Bbb R_+$ будем 


обозначать множество всех неотрицательных действительных чисел. 


 


\section{Условия на операторные коэффициенты \protect\\


и данные рассматриваемой задачи} 


 


Для вещественного сепарабельного гильбертова пространства $H$, в котором 


рассматривается задача (1.1), (1.2), обозначим через 


$(\cdot,\cdot{}\cdot)_H$ и $|\cdot|_H$ скалярное 


произведение и норму этого пространства соответственно. Все операторы 


$K(t)$ ($t\in[0,T]$), участвующие в уравнении (1.1), принадлежат 


$\cal L(H)$ и удовлетворяют следующему условию. 


\medskip


 


{\it Условие} (K). Функция $t\to K(t)$ имеет на $[0,T]$ обобщенную 


производную $t\to K'(t)$ (из $L(0,T;\cal L(H))$\,). 


\medskip 


 


В дальнейшем будет использоваться обозначение 


$$ 


k_0=\max_{0\le t\le T}\|K(t)\mid\cal L(H)\|. 


$$ 


 


Все операторы $A(t)$ ($t\in[0,T]$), участвующие в (1.1), определены на 


одном и том же линейном множестве $D(A)$, плотном в $H$, и принимают 


значения в $H$. Все они 


неограниченные, самосопряженные и положительно-определенные. Обозначим 


через $a(t,\cdot,\cdot{}\cdot)$ скалярное произведение энергетического 


пространства (\cite{1}, 


с.\,21), порожденного оператором $A(t)$. Далее предположим, что 


существует 


самосопряженный положительно-определенный оператор $A_0$, действующий в 


$H$, определенный на $D(A)$ (т.\,е.~сходный с каждым из операторов 


$A(t)$\,)\footnote{$^)$Поскольку операторы $A_0$ и $A(t)$ сходные, все 


операторы $A(t)A^{-1}_0$, 


$A_0(A(t))^{-1}$, $(A(t))^{1/2}A^{-1/2}_0$, $A^{1/2}_0(A(t))^{-1/2}$ 


ограничены в $H$ (\cite{1}, с.\,22--23).}{$^)$}, имеющий 


вполне непрерывный в $H$ обратный оператор $A^{-1}_0$ и образующий с 


каждым из 


операторов $A(t)$ острый угол. Энергетическое пространство, порожденное 


оператором $A_0$, его скалярное произведение и норму обозначим 


соответственно через $H_0$, $(\cdot,\cdot{}\cdot)_0$, $|\cdot|_0$. 


Предполагается, 


что операторы $A(t)$ ($t\in[0,T]$) и $A_0$ удовлетворяют следующим 


четырем условиям. 


 


{\it Условие} (A1). Существует константа $a_0>0$ такая, что для всех 


$t\in[0,T]$


$$


\|A(t)A^{-1}_0\mid\cal L(H)\|\le a_0. 


$$


 


{\it Условие} (A2). Существует константа $a_1>0$ такая, что для всех 


$t\in[0,T]$


$$


\|A^{1/2}_0(A(t))^{-1/2}\mid\cal L(H)\|\le a_1. 


$$


 


{\it Условие} (A3). Существует константа $a_2>0$ такая, что для всех 


$t\in[0,T]$, $u\in D(A)$


$$


(A(t)u,A_0u)_H\ge a_2|A_0u|^2_H. 


$$


 


{\it Условие} (A4). Функция $t\to a(t,\cdot,\cdot{}\cdot)$ принадлежит 


$C^1([0,T];\cal B(H_0))$ и имеет на $[0,T]$ обобщенную производную 


второго порядка (из $L(0,T;\cal B(H_0))$\,). 


 


Ниже будут использоваться обозначения: $t\to a'(t,\cdot,\cdot{}\cdot)$ 


--- производная первого порядка функции $t\to a(t,\cdot,\cdot{}\cdot)$; 


$$ 


\alpha_0=\max_{0\le t\le T}\|a'(t,\cdot,\cdot{}\cdot)\mid 


\cal B(H_0)\|. 


$$ 


 


Нелинейный оператор $B$ действует из $[0,T]\times H_0$ в $H$ и 


удовлетворяет следующему условию. 


 


{\it Условие} (B). В каждой точке $(t,u)\in[0,T]\times H_0$ существует 


производная Фреше $B'[t,u]\in\cal L(\Bbb R\oplus H_0,H)$ 


оператора $B$, причем существуют функции 


$\beta_i:[0,T]\times\Bbb R_+\to\Bbb R_+$ ($i=0,1$) такие, что для почти 


всех (п.\,в.) $t\in[0,T]$ функции $\beta_i(t,\cdot)$ ($i=0,1$) 


неубывающие на $\Bbb R_+$, для любого 


$\xi\in\Bbb R_+$\; $\beta_i(\cdot,\xi)\in L(0,T)$ ($i=0,1$) и 


$|B'[t,u](\tau,v)|_H\le\beta_0(t,|u|_0)|\tau|+ 


\beta_1(t,|u|_0)|v|_0$ для всех 


$(t,u)\in[0,T]\times H_0$, $(\tau,v)\in\Bbb R\times H_0$. 


 


Кроме того, предполагается, что существует некоторый функционал 


$\Phi:[0,T]\times D(A)\to\Bbb R_+$, 


удовлетворяющий следующим условиям. 


 


{\it Условие} ($\Phi 1$). Существует неубывающая функция 


$\varphi_0:\Bbb R_+\to\Bbb R_+$ такая, что для всех 


$u\in D(A)$


$$


\Phi(0,u)\le\varphi_0(|A_0u|_H).


$$


 


{\it Условие} ($\Phi 2$). Существует неубывающая функция 


$\varphi:\Bbb R_+\to\Bbb R_+$ такая, что 


для всех $(t,u)\in[0,T]\times D(A)$\; $|u|_0\le\varphi(\Phi(t,u))$. 


 


В условии $(\Phi 3)$ подразумевается, что на $D(A)$ введено скалярное 


произведение 


$(\cdot,\cdot{}\cdot)_{D(A)}=(A_0\cdot,A_0\cdot{}\cdot)_H$. 


 


{\it Условие} ($\Phi 3$). В каждой точке $(t,u)\in[0,T]\times D(A)$ 


существует производная 


Фреше $\Phi'[t,u]\in\cal L(\Bbb R\oplus D(A),\Bbb R)$ функционала 


$\Phi$, причем существуют 


функции $\varphi_i:[0,T]\to\Bbb R_+$ ($i=1,2,3$), принадлежащие 


$L(0,T)$ и такие, что для всех 


$(t,u)\in[0,T]\times D(A)$, $v\in D(A)$


$$


\Phi'[t,u](1,v)-a(t,u,v)-(B(t,u),v)_H\le\varphi_1(t)+ 


\varphi_2(t)|v|^2_H+\varphi_3(t)\Phi(t,u). 


$$


 


На начальные данные (1.2) наложим условия 


$u^{(0)}_0\in D(A)$, $u^{(0)}_1\in H_0$. Все условия (K), 


(A1)--(A4), (B), ($\Phi 1$)--($\Phi 3$) и последние условия на 


начальные данные считаются выполненными всюду на протяжении данной 


статьи. 


 


\begin{notenonum} 


Важным для приложений частным случаем уравнения (1.1) является 


уравнение с нелинейностью вида 


$B(t,u(t))=B_0(u(t))+f(t)$, где $B_0:H_0\to H$ --- потенциальный 


оператор, 


функция $f$ принадлежит, например, $C^1([0,T];H)$. Для таких уравнений в 


качестве $\Phi$ 


в ряде случаев подходит функционал 


$(t,u)\to(1/2)a(t,u,u)+\Phi_0(u)$, где $\Phi_0$ --- потенциал оператора 


$B_0$. Задача (1.1), (1.2) с нелинейностью указанного вида и не 


зависящими 


от $t$ операторами $K$, $A$ может интерпретироваться как абстрактная 


формулировка начально-краевых задач математической физики, получаемых 


на основе вариационного принципа Гамильтона--Остроградского (см. 


\cite{10}, с.\,749--751). Задачи этого класса, к которым 


непосредственно применимы результаты данной работы, рассматривались в 


(\cite{12}, гл.\,I, \S\,1; \cite{13}). Другим частным случаем уравнения 


(1.1) является линейное уравнение с 


$B(t,u(t))=B_1(t)u(t)+f(t)$, где, например, 


$B_1\in C^1([0,T];\cal L(H_0,H))$, $f\in C^1([0,T];H)$. В 


этом случае следует принять $\Phi(t,u)=(1/2)a(t,u,u)$. 


\end{notenonum} 


 


Отметим необходимые для дальнейшего следствия из условия (B). 


 


\begin{lem} 


Для п.\,в.{} $t\in[0,T]$ $($именно, для всех $t$, при которых функция 


$\beta_1(t,\cdot)$ не убывает на $\Bbb R_+)$ и для любых $u,v\in H_0$ 


\begin{equation} 


|B(t,u)-B(t,v)|_H\le \beta_1(t,\max(|u|_0,|v|_0))|u-v|_0. 


\end{equation} 


\end{lem} 


 


\begin{pf} 


Из условия (B) следует, что для любого $t\in[0,T]$ оператор 


$B_t:H_0\to H$, 


определенный 


равенством $B_t(u)=B(t,u)$, имеет в каждой точке $u\in H_0$ производную 


Фреше $B'_t[u]\in\cal L(H_0,H)$, причем 


$\|B'_t[u]\mid\cal L(H_0,H)\|\le\beta_1(t,|u|_0)$ 


для всех $u\in H_0$. Поэтому, используя формулу конечных приращений 


(\cite{14}, 


с.\,483), получаем для любых $u,v\in H_0$ и для всех $t\in[0,T]$, при 


которых функция $\beta_1(t,\cdot)$ не убывает на $\Bbb R_+$, 


$|B(t,u)-B(t,v)|_H\le\sup\limits_{0\le\theta\le 1} 


\beta_1(t,|\theta u+(1-\theta)v|_0)|u-v|_0= 


\beta_1(t,\max(|u|_0,|v|_0))|u-v|_0$. 


\end{pf} 


 


\begin{lem} 


Существует неубывающая функция $\beta:\Bbb R_+\to\Bbb R_+$ такая, что 


$|B(t,u)|_H\le\beta(|u|_0)$ для всех $(t,u)\in[0,T]\times H_0$. 


\end{lem} 


 


\begin{pf} 


Пусть $t_0\in[0,T]$ --- такое число, что функция $\beta_1(t_0,\cdot)$ 


не убывает на $\Bbb R_+$. Зафиксируем 


произвольные $t\in[0,T]$, $u\in H_0$. 


Из условия (B) следует, что функция $B_u:[0,T]\to H$, определенная 


равенством $B_u(\tau)=B(\tau,u)$, дифференцируема в каждой 


точке $\tau\in[0,T]$ и 


для всех $\tau\in[0,T]$ выполняется 


$\big|\frac{d}{d\tau}B_u(\tau)\big|_H=|B'[\tau,u](1,0)|_H\le 


\beta_0(\tau,|u|_0)$. Отсюда с 


помощью теоремы 1 из (\cite{15}, с.\,251) получим 


\begin{equation} 


|B(t,u)-B(t_0,u)|_H\le 


\int^{\max(t_0,t)}_{\min(t_0,t)}\beta_0(\tau,|u|_0)d\tau. 


\end{equation} 


Используя оценки (2.1) и (2.2), имеем 


$|B(t,u)|_H\le|B(t,u)-B(t_0,u)|_H+|B(t_0,u)-B(t_0,0)|_H+ 


|B(t_0,0)|_H\le\beta(|u|_0)$, 


где $\beta:\Bbb R_+\to\Bbb R_+$ --- функция, 


определенная равенством 


$$ 


\beta(\xi)=\max\bigg\{\,\int^{t_0}_0\beta_0(\tau,\xi)d\tau,\; 


\int^T_{t_0}\beta_0(\tau,\xi)d\tau\bigg\}+\beta_1(t_0,\xi)\xi+ 


|B(t_0,0)|_H. 


$$ 


Очевидно, эта функция не убывает на $\Bbb R_+$. 


\end{pf} 


 


В заключение этого пункта сформулируем определение решения 


рассматриваемой задачи, действующее всюду ниже в данной работе. 


 


\begin{defnnonum} 


Функцию $u_*:[0,T]\to D(A)$ будем называть решением задачи (1.1), (1.2), 


если $A_0u_*\in L^\infty(0,T;H)$, $u'_*\in L^\infty(0,T;H_0)$, 


$u''_*\in L^\infty(0,T;H)$, для 


п.\,в.{} $t\in[0,T]$\; $Q[u_*,t]=0$, $u_*(0)=u^{(0)}_0$ и 


$u'_*(0)=u^{(0)}_1$. 


\end{defnnonum} 


 


\section{Специальный базис и соотношения МБГ. Существование, \protect\\


единственность и априорные оценки приближенных решений} 


 


Поскольку оператор $A_0$ имеет вполне непрерывный в $H$ обратный 


оператор $A^{-1}_0$, то в силу теоремы Гильберта--Шмидта (\cite{14}, 


с.\,246) в $H$ 


существует базис $\{v_{0,m}\}^\infty_{m=1}$, состоящий из собственных 


элементов оператора $A_0$, 


которые можно считать пронумерованными в порядке неубывания


соответствующих собственных 


значений $\lambda_m$ ($m=1,2,\dotsc$), $A_0v_{0,m}=\lambda_m v_{0,m}$, 


$0<\lambda_m\le\lambda_{m+1}$, $m=1,2,\dotsc$; 


$\lim\limits_{m\to\infty}\lambda_m=\infty$. Пусть $n$ --- произвольное 


натуральное число. Обозначим через $H_n$ линейную оболочку конечной 


системы элементов 


$\{v_{0,m}\}^n_{m=1}$, через $P_n$ --- ортопроектор $H$ на $H_n$ и 


через $R_n$ --- оператор $I-P_n$, где $I$ --- тождественный оператор в 


$H$. Будем решать задачу (1.1), (1.2) по 


МБГ, используя в качестве базисных элементов собственные элементы 


оператора $A_0$


\begin{align} 


& P_n(Q[u_n,t])=0,\quad t\in[0,T], \\ 


& u_n(t)\in H_n,\quad t\in [0,T], \\ 


& u_n(0)=P_n u^{(0)}_0,\quad u'_n(0)=P_n u^{(0)}_1. 


\end{align} 


Решение $u_n$ задачи (3.1)--(3.3) будем называть приближенным решением 


задачи (1.1), (1.2), построенным по МБГ. 


 


В дальнейшем через $c_i$ ($i=1,2,\dotsc,8$) обозначаются различные 


неотрицательные 


константы, которые могут зависеть только от вида операторов $K(t)$, 


$A(t)$ ($t\in[0,T]$), $A_0$, $B$, функционала $\Phi$ и от $u^{(0)}_0$, 


$u^{(0)}_1$, $T$. Восстанавливая 


описываемые ниже построения во всех деталях, можно получить значения 


всех этих констант в явном виде. Через $\|\cdot\|_C$ в дальнейшем 


обозначается следующая норма пространства 


$C([0,T];H_0)\cap C^1([0,T];H)$:


$$ 


\|u\|_C=\max_{0\le t\le T} 


(|u'(t)|^2_H+|u(t)|^2_0)^{1/2}. 


$$ 


 


Следуя (\cite{10}, \cite{12}, сс.\,23--24, 30--33), нетрудно показать, 


что при любом натуральном $n$ задача (3.1)--(3.3) имеет единственное 


решение $u_n$, удовлетворяющее условиям 


$u_n\in C^2([0,T];H_n)$, $u'''_n\in L(0,T;H_n)$, и для этого решения 


выполнены априорные оценки 


\begin{align} 


\|u_n\|_C &\le c_1, \\ 


\|u'_n\|_C &\le c_2. 


\end{align} 


 


\begin{notenonum} 


Необходимым элементом вывода оценок (3.4), (3.5) и нижеследующих 


построений является лемма Гронуолла, используемая в данной работе в 


следующей формулировке, усиленной по сравнению с более известными 


формулировками из (\cite{16}, с.\,298; \cite{17}, с.\,37): если 


$y\in C([0,T])$, $y_0\in\Bbb R$, 


$z\in L(0,T)$ и для всех $t\in[0,T]$\; $z(t)\ge 0$, 


$y(t)\le{}y_0+\int\limits^t_0{}y(\tau)z(\tau)d\tau$, то 


$y(t)\le y_0\exp\Big(\int\limits^t_0z(\tau)d\tau\Big)$ 


для всех $t\in[0,T]$. Также отметим, что функционал $\Phi$ используется 


в данной работе только при выводе оценки (3.4). 


\end{notenonum} 


 


\section{Основные результаты} 


 


\begin{thmnonum} 


Решение $u_*$ задачи $(1.1)$, $(1.2)$ существует и единственно, 


последовательность $\{u_n\}^\infty_{n=1}$ приближенных решений задачи 


$(1.1)$, $(1.2)$, построенных по МБГ $(3.1)$--$(3.3)$, сильно сходится 


к $u_*$ в 


пространстве $C([0,T];H_0)\cap C^1([0,T];H)$ 


и скорость этой сходимости допускает оценку 


\begin{equation} 


\|u_n-u_*\|_C\le c_3|R_n u^{(0)}_1|_H+c_4 


|R_n u^{(0)}_0|_0+c_5\lambda^{-1/4}_{n+1},


\end{equation} 


где $n$ --- произвольное натуральное число. 


\end{thmnonum} 


 


\begin{pf} 


Прежде всего установим априорную оценку нормы 


$\|A_0u_n\mid C([0,T];H)\|$ ($n$ --- любое 


натуральное число, $u_n$ --- решение задачи (3.1)--(3.3)\,). Умножим 


скалярно в $H$ обе части равенства (3.1) на $A_0u_n(t)$. Тогда получим 


$(A(t)u_n(t),A_0u_n(t))_H=-(u''_n(t)+K(t)u'_n(t)+B(t,u_n(t)), 


A_0u_n(t))_H$ для всех $t\in[0,T]$. Оценим левую часть данного 


равенства снизу с помощью условия 


(A3), а правую часть --- сверху через 


$[|u''_n(t)|_H+|K(t)u'_n(t)|_H+|B(t,u_n(t))|_H]\,|A_0u_n(t)|_H$. 


Таким образом, получим неравенство 


$|A_0u_n(t)|_H\le a^{-1}_2[|u''_n(t)|_H+|K(t)u'_n(t)|_H+ 


|B(t,u_n(t))|_H]$, 


правую часть которого оценим с помощью (3.4), (3.5), 


очевидного неравенства $\|K(t)\mid\cal L(H)\|\le k_0$ и леммы 2. В 


результате имеем 


\begin{equation} 


\|A_0u_n\mid C([0,T];H)\|\le c_6. 


\end{equation} 


Далее установим априорную оценку нормы $|Q[u_n,t]|_H$, где $t$ 


--- любое число из $[0,T]$. В силу условия (A1)\; 


$|Q[u_n,t]|_H\le|u''_n(t)|_H+|K(t)u'_n(t)|_H+a_0|A_0u_n(t)|_H+ 


|B(t,u_n(t))|_H$. Используя это неравенство, (3.4), 


(3.5), (4.2), неравенство $\|K(t)\mid\cal L(H)\|\le k_0$ 


и лемму 2, имеем 


\begin{equation} 


|Q[u_n,t]|_H\le c_7. 


\end{equation} 


 


Теперь покажем, что последовательность $\{u_n\}^\infty_{n=1}$ 


приближенных решений задачи 


(1.1), (1.2), построенных по МБГ (3.1)--(3.3), сильно сходится в 


пространстве $C([0,T];H_0)\cap C^1([0,T];H)$. Пусть $n$, $s$ --- 


произвольные натуральные числа такие, 


что $n\le s$. Примем обозначение $\Delta u_{n,s}=u_n-u_s$. Из (3.1) 


следует, что для всех $t\in[0,T]$\; 


$P_s(Q[u_n,t])=R_n P_s(Q[u_n,t])$. 


Вычтем из этого равенства равенство (3.1) с $s$ вместо $n$, обе части 


полученного равенства умножим скалярно в $H$ на 


$2\Delta u'_{n,s}(t)$ и вновь полученное равенство почленно 


проинтегрируем по отрезку $[0,\tau]$, где $\tau$ --- 


произвольное число из $[0,T]$, с учетом начальных условий 


$\Delta u_{n,s}(0)=P_n u^{(0)}_0-P_s u^{(0)}_0$, 


$\Delta u'_{n,s}(0)=P_n u^{(0)}_1-P_s u^{(0)}_1$. В 


результате будем иметь для всех $\tau\in[0,T]$ 


\begin{multline} 


|\Delta u'_{n,s}(\tau)|^2_H+a(\tau,\Delta u_{n,s}(\tau), 


\Delta u_{n,s}(\tau))=|P_n u^{(0)}_1-P_s u^{(0)}_1|^2_H+ 


a(0,P_n u^{(0)}_0-P_s u^{(0)}_0,P_n u^{(0)}_0- 


P_s u^{(0)}_0) - \\ 


% 


-2\int^\tau_0(K(t)\Delta u'_{n,s}(t),\Delta u'_{n,s}(t))_H dt+ 


2\int^\tau_0(B(t,u_s(t))-B(t,u_n(t)),\Delta u'_{n,s}(t))_H dt +\\ 


% 


+2\int^\tau_0(R_n P_s(Q[u_n,t]),\Delta u'_{n,s}(t))_H dt+ 


\int^\tau_0 a'(t,\Delta u_{n,s}(t),\Delta u_{n,s}(t))dt. 


\end{multline} 


В силу условия (A2) второе слагаемое в левой части (4.4) допускает 


оценку снизу 


\begin{equation} 


a(\tau,\Delta u_{n,s}(\tau),\Delta u_{n,s}(\tau))\ge 


a^{-2}_1|\Delta u_{n,s}(\tau)|^2_0. 


\end{equation} 


Оценим сверху каждое слагаемое в правой части (4.4) в отдельности. Для 


первого слагаемого получим 


\begin{equation} 


|P_n u^{(0)}_1-P_s u^{(0)}_1|^2_H=|P_s R_n u^{(0)}_1|^2_H\le 


|R_n u^{(0)}_1|^2_H. 


\end{equation} 


Для второго слагаемого имеем 


\begin{multline} 


a(0,P_n u^{(0)}_0-P_s u^{(0)}_0, P_n u^{(0)}_0-P_s u^{(0)}_0)= 


|(A(0))^{1/2}P_s R_n u^{(0)}_0|^2_H\le \\ 


% 


\le \|(A(0))^{1/2}A^{-1/2}_0\mid\cal L(H)\|^2 


|A^{1/2}_0 P_s R_n u^{(0)}_0|^2_H \le \\ 


% 


\le \|(A(0))^{1/2}A^{-1/2}_0\mid\cal L(H)\|^2 


|R_n u^{(0)}_0|^2_0. 


\end{multline} 


Третье слагаемое оценивается следующим образом:


\begin{multline} 


-2\int^\tau_0(K(t)\Delta u'_{n,s}(t),\Delta u'_{n,s}(t))_H dt \le \\


% 


\le 2\int^\tau_0\|K(t)\mid\cal L(H)\|\, 


|\Delta u'_{n,s}(t)|^2_H dt \le 2k_0\int^\tau_0 


|\Delta u'_{n,s}(t)|^2_H dt. 


\end{multline} 


Используя лемму 1, оценку (3.4) и неравенство Коши, получим 


\begin{multline} 


2\int^\tau_0(B(t,u_s(t))-B(t,u_n(t)), 


\Delta u'_{n,s}(t))_H dt \le 2\int^\tau_0 \beta_1(t,c_1) 


|\Delta u_{n,s}(t)|_0|\Delta u'_{n,s}(t)|_H dt \le \\ 


% 


\le\int^\tau_0\beta_1(t,c_1)[|\Delta u_{n,s}(t)|^2_0+ 


|\Delta u'_{n,s}(t)|^2_H]dt. 


\end{multline} 


В силу специального выбора базиса МБГ\; 


$\|A^{-1/2}_0R_n\mid\cal L(H)\|=\lambda^{-1/2}_{n+1}$, в силу (3.5) для 


всех $t\in[0,\tau]$\; 


$|\Delta u'_{n,s}(t)|_0\le 2c_2$ и в соответствии с (4.3) для всех 


$t\in[0,\tau]$\; $|Q[u_n,t]|_H\le c_7$, поэтому 


\begin{multline} 


2\int^\tau_0(R_n P_s(Q[u_n,t]),\Delta u'_{n,s}(t))_H dt= 


2\int^\tau_0(A^{-1/2}_0R_n P_s(Q[u_n,t]), 


A^{1/2}_0\Delta u'_{n,s}(t))_H dt \le \\ 


% 


\le 2\|A^{-1/2}_0R_n\mid\cal L(H)\|\int^\tau_0 


|Q[u_n,t]|_H\,|\Delta u'_{n,s}(t)|_0 dt\le 


4T c_2c_7\lambda^{-1/2}_{n+1}. 


\end{multline} 


Наконец, последнее слагаемое в правой части (4.4) допускает оценку 


\begin{equation} 


\int\limits^\tau_0 a'(t,\Delta u_{n,s}(t),\Delta u_{n,s}(t))dt\le 


\int\limits^\tau_0\|a'(t,\cdot,\cdot{}\cdot)\mid\cal B(H_0)\|\, 


|\Delta u_{n,s}(t)|^2_0dt \le  


\alpha_0\int\limits^\tau_0|\Delta u_{n,s}(t)|^2_0dt. 


\end{equation} 


Из (4.4)--(4.11) следует, что для всех $\tau\in[0,T]$ 


\begin{multline*} 


|\Delta u'_{n,s}(\tau)|^2_H+a^{-2}_1|\Delta u_{n,s}(\tau)|^2_0\le 


|R_n u^{(0)}_1|^2_H+\|(A(0))^{1/2}A^{-1/2}_0\mid\cal L(H)\|^2\, 


|R_n u^{(0)}_0|^2_0 + \\ 


% 


+ 4T c_2c_7\lambda^{-1/2}_{n+1}+\int^\tau_0 


[(\beta_1(t,c_1)+2k_0)|\Delta u'_{n,s}(t)|^2_H+ 


(\beta_1(t,c_1)+\alpha_0)|\Delta u_{n,s}(t)|^2_0]dt, 


\end{multline*} 


откуда с помощью леммы Гронуолла получим 


\begin{equation} 


\|u_n-u_s\|_C\le c_3|R_n u^{(0)}_1|_H+c_4|R_n u^{(0)}_0|_0+


c_5\lambda^{-1/4}_{n+1}. 


\end{equation} 


Из (4.12) видно, что последовательность $\{u_n\}^\infty_{n=1}$ 


фундаментальна в 


пространстве $C([0,T];H_0)\cap C^1([0,T];H)$. 


Следовательно, существует функция 


$u_*\in C([0,T];H_0)\cap C^1([0,T];H)$ такая, что 


\begin{equation} 


u_n\to u_*\quad 


\text{сильно в}\quad 


C([0,T];H_0)\cap C^1([0,T];H),\quad n\to\infty. 


\end{equation} 


 


Теперь покажем, что функция $u_*$ является решением задачи (1.1), (1.2). 


Примем обозначения: 


$\la v_1,v_2,t\ra=(v''_1(t),v''_2(t))_H+(v'_1(t),v'_2(t))_0+ 


(A_0v_1(t),A_0v_2(t))_H$, 


где $t\in[0,T]$, $v_1$ и $v_2$ --- произвольные функции, 


отображающие $[0,T]$ в $D(A)$ и удовлетворяющие условиям 


$A_0v_i\in L^2(0,T;H)$, $v'_i\in L^2(0,T;H_0)$, $v''_i\in L^2(0,T;H)$ 


($i=1,2$); $\cal H$ 


--- сепарабельное гильбертово пространство, состоящее из функций 


$v:[0,T]\to D(A)$, 


удовлетворяющих условиям 


$A_0v\in L^2(0,T;H)$, $v'\in L^2(0,T;H_0)$, $v''\in L^2(0,T;H)$, и 


снабженное скалярным произведением 


$(v_1,v_2)_{\cal H}=\int\limits^T_0\la v_1,v_2,t\ra dt$. В силу (3.5), 


(4.2) для всех $n\in\Bbb N$ 


\begin{equation} 


\max_{0\le t\le T}\la u_n,u_n,t\ra\le c_8. 


\end{equation} 


Тем более последовательность $\{u_n\}^\infty_{n=1}$ ограничена в 


$\cal H$. Следовательно, по 


теореме 3 из (\cite{14}, с.\,199) найдется подпоследовательность 


$\{u_{n_k}\}^\infty_{k=1}$ этой последовательности, слабо сходящаяся в 


$\cal H$ к некоторой функции $v_*$. 


Поскольку эта подпоследовательность к тому же сильно сходится в 


$C([0,T];H_0)\cap C^1([0,T];H)$ к 


$u_*$, то в силу единственности слабого предела в 


$C([0,T];H_0)\cap C^1([0,T];H)$\; $v_*=u_*$ и вся 


последовательность $\{u_n\}^\infty_{n=1}$ слабо сходится в $\cal H$ к 


$u_*$. 


 


Сейчас будет доказано, что функция $u_*$ удовлетворяет условиям 


(участвующим в определении решения задачи (1.1), (1.2), см.{} п.\,2) 


\begin{equation} 


A_0u_*\in L^\infty(0,T;H),\quad 


u'_*\in L^\infty(0,T;H_0),\quad 


u''_*\in L^\infty(0,T;H). 


\end{equation} 


Пусть $m_0$ --- какое-либо фиксированное число, строго большее, чем 


$c_8$. Обозначим через $\sigma$ множество 


$\{t\in[0,T]\mid\la u_*,u_*,t\ra\ge m_0\}$. Рассмотрим линейный 


ограниченный функционал $F$ над $\cal H$, определенный равенством 


$F(v)=\int\limits_\sigma\la u_*,u_*,t\ra^{-1/2}\la v,u_*,t\ra dt$. В 


силу (4.14) для всех $n\in\Bbb N$ 


\begin{equation} 


F(u_n)\le\int_\sigma\la u_n,u_n,t\ra^{1/2}dt\le 


c^{1/2}_8\mes\sigma. 


\end{equation} 


Так как последовательность $\{u_n\}^\infty_{n=1}$ слабо сходится в 


$\cal H$ к $u_*$, то 


\begin{equation} 


F(u_*)=\lim_{n\to\infty}F(u_n). 


\end{equation} 


Из (4.16) и (4.17) получаем $F(u_*)\le c^{1/2}_8\mes\sigma$. С другой 


стороны,


$F(u_*)=\int\limits_\sigma\la u_*,u_*,t\ra^{1/2}dt\ge m^{1/2}_0 


\mes\sigma$ по определению множества $\sigma$. Из двух последних 


неравенств вытекает $m^{1/2}_0\mes\sigma\le c^{1/2}_8\mes\sigma$, но при 


этом $m_0>c_8\ge 0$. 


Очевидно, это возможно только при $\mes\sigma=0$. Поэтому 


$\la u_*,u_*,t\ra<m_0$ для п.\,в.{} $t\in[0,T]$ и, стало быть, функция 


$u_*$ удовлетворяет условиям (4.15). 


 


Из того, что $u_n\to u_*$ слабо в $\cal H$, следует 


\begin{alignat}{2} 


A_0u_n &\to A_0u_* &\ \; &\text{слабо в \;}L^2(0,T;H), 


\quad n\to \infty, \\ 


u''_n &\to u''_* &\ \; &\text{слабо в \;}L^2(0,T;H), 


\quad n\to \infty. 


\end{alignat} 


Далее получим из (4.13) с учетом условия (K) 


\begin{equation} 


K(\cdot)u'_n(\cdot)\to K(\cdot)u'_*(\cdot)\quad 


\text{сильно в}\quad 


C([0,T];H),\quad n\to \infty, 


\end{equation} 


а из (4.18) с учетом условия (A1) 


\begin{equation} 


A(\cdot)u_n(\cdot)\to A(\cdot)u_*(\cdot)\quad 


\text{слабо в}\quad 


L^2(0,T;H),\quad n\to\infty. 


\end{equation} 


В силу (4.13) в (3.4) можно перейти к пределу. Используя 


устанавливаемую таким образом оценку $\|u_*\|_C\le c_1$, саму оценку 


(3.4) и лемму 1, имеем для п.\,в.{} $t\in[0,T]$\; 


$|B(t,u_n(t))-B(t,u_*(t))|_H\le\beta_1(t,c_1)|u_n(t)-u_*(t)|_0$, 


поэтому 


\begin{equation} 


B(\cdot,u_n(\cdot))\to B(\cdot,u_*(\cdot))\quad 


\text{сильно в}\quad 


L(0,T;H),\quad n\to\infty. 


\end{equation} 


Теперь, используя (4.13), (4.19)--(4.22), предельным переходом в (3.1), 


(3.3) получаем равенства $Q[u_*,t]=0$ (для п.\,в.{} 


$t\in[0,T]$), $u_*(0)=u^{(0)}_0$, $u'_*(0)=u^{(0)}_1$. В 


частности, 


первое из этих равенств устанавливается следующим образом: полагая 


$m,n\in\Bbb N$, 


$n\ge m$, умножая скалярно (3.1) в $H$ на 


$\text{sign\,}((Q[u_*,t],v_{0,m})_H)v_{0,m}$, интегрируя по $[0,T]$ и 


переходя к 


пределу при $n\to\infty$, получаем 


$\int\limits^T_0|(Q[u_*,t],v_{0,m})_H|dt=0$ 


для любого $m\in\Bbb N$, т.\,е.{} 


$(Q[u_*,t],v_{0,m})_H=0$ 


для п.\,в.{} $t\in[0,T]$, и 


затем из последнего равенства (поскольку $\{v_{0,m}\}^\infty_{m=1}$ --- 


базис в $H$) заключаем, что $Q[u_*,t]=0$ для п.\,в.{} $t\in[0,T]$. 


 


Итак, сильный предел $u_*$ последовательности $\{u_n\}^\infty_{n=1}$ в 


пространстве $C([0,T];H_0)\cap C^1([0,T];H)$ 


удовлетворяет всем условиям определения решения задачи (1.1), (1.2), 


так что утверждения теоремы о существовании решения этой задачи и 


сходимости к нему последовательности приближенных решений доказаны. 


Переходя в (4.12) к пределу при $s\to\infty$, имеем оценку (4.1). 


 


Остается доказать, что решение задачи (1.1), (1.2) единственно. 


Следующее доказательство аналогично проведенному в (\cite{12}, 


с.\,27--28) для конкретной начально-краевой задачи. Пусть $u^{(1)}_*$, 


$u^{(2)}_*$ --- два решения задачи (1.1), (1.2) и 


$\Delta u=u^{(1)}_*-u^{(2)}_*$, тогда для 


п.\,в.{} $t\in[0,T]$ 


\begin{gather} 


\Delta u''(t)+K(t)\Delta u'(t)+A(t)\Delta u(t)+ 


B(t,u^{(1)}_*(t))-B(t,u^{(2)}_*(t))=0,\\ 


\Delta u(0)=\Delta u'(0)=0. 


\end{gather} 


Умножим обе части равенства (4.23) скалярно в $H$ на 


$2\Delta u'(t)$ и полученное 


равенство почленно проинтегрируем по отрезку $[0,\tau]$, где $\tau$ --- 


произвольное число из $[0,T]$, с учетом начальных условий (4.24). 


Тогда будем иметь для всех $\tau\in[0,T]$ 


\begin{multline*} 


|\Delta u'(\tau)|^2_H+a(\tau,\Delta u(\tau),\Delta u(\tau))= \\ 


% 


=\int^\tau_0[2(-K(t)\Delta u'(t)+B(t,u^{(2)}_*(t))- 


B(t,u^{(1)}_*(t)),\Delta u'(t))_H+ 


a'(t,\Delta u(t),\Delta u(t))]dt. 


\end{multline*} 


Оценивая левую часть данного равенства снизу с помощью условия (A2), а 


правую --- сверху с помощью очевидных неравенств 


$\|K(t)\mid\cal L(H)\|\le k_0$, 


$\|a'(t,\cdot,\cdot{}\cdot)\mid\cal B(H_0)\|\le\alpha_0$, леммы 1 и 


неравенства Коши, получим для всех $\tau\in[0,T]$ 


$$


|\Delta u'(\tau)|^2_H+a^{-2}_1|\Delta u(\tau)|^2_0\le  


\int^\tau_0[(\beta_1(t,M)+2k_0) 


|\Delta u'(t)|^2_H+(\beta_1(t,M)+\alpha_0) 


|\Delta u(t)|^2_0]dt, 


$$


где $M=\max(\|u^{(1)}_*\mid C([0,T];H_0)\|,\, 


\|u^{(2)}_*\mid C([0,T];H_0)\|)$. 


Из последнего неравенства с помощью леммы Гронуолла получаем 


$\Delta u(\tau)=0$ 


для всех $\tau\in[0,T]$, т.\,е.{} $u^{(1)}_*=u^{(2)}_*$. 


Таким образом, решение задачи (1.1), (1.2) единственно. 


\end{pf} 


 


\begin{thebibliography}{99} 


 


\bibitem{1} 


Михлин С.Г. {\em Численная реализация вариационных методов}. -- М.: 


Наука, 1966. -- 432 с. 


\bibitem{2} 


Соболевский П.Е. {\em Об уравнениях с операторами, образующими острый 


угол} // ДАН СССР. -- 1957. -- Т.\,116. -- \No\,5. -- С.\,752--757. 


\bibitem{3} 


Железовский С.Е., Кириченко В.Ф., Крысько В.А. {\em О скорости 


сходимости метода Бубнова--Гал\"еркина для одной неклассической системы 


дифференциальных уравнений} // Дифференц. уравнения. -- 1987. -- 


Т.\,23. -- \No\,8. -- С.\,1407--1416. 


\bibitem{4} 


Железовская Л.А., Железовский С.Е., Кириченко В.Ф., Крысько В.А. {\em О 


скорости сходимости метода Бубнова--Гал\"еркина для гиперболических 


уравнений} // Дифференц. уравнения. -- 1990. -- Т.\,26. -- \No\,2. -- 


С.\,323--333. 


\bibitem{5} 


Железовский С.Е. {\em О скорости сходимости метода Бубнова--Гал\"еркина 


для нелинейной задачи теории упругих оболочек} // Дифференц. уравнения. 


-- 1995. -- Т.\,31. -- \No\,5. -- С.\,858--869. 


\bibitem{6} 


Джишкариани А.В. {\em О быстроте сходимости метода Бубнова--Гал\"еркина} 


// Журн. вычисл. матем. и матем. физ. -- 1964. -- Т.\,4. -- \No\,2. -- 


С.\,343--348. 


\bibitem{7} 


Вайникко Г.М. {\em Некоторые оценки погрешности метода 


Бубнова--Гал\"еркина}. 2. {\em Оценки $n$-го приближения} // Учен. зап. 


Тартуск. ун-та. -- 1964. -- \No\,150. -- С.\,202--215. 


\bibitem{8} 


Зарубин А.Г. {\em Исследование проекционной процедуры 


Гал\"еркина--Петрова методом дробных степеней} // ДАН СССР. -- 1987. -- 


Т.\,297. -- \No\,4. -- С.\,780--784. 


\bibitem{9} 


Зарубин А.Г. {\em О скорости сходимости метода Фаэдо--Гал\"еркина для 


квазилинейных нестационарных операторных уравнений} // Дифференц. 


уравнения. -- 1990. -- Т.\,26. -- \No\,12. -- С.\,2051--2059. 


\bibitem{10} 


Ворович И.И. {\em О некоторых прямых методах в нелинейной теории 


колебаний пологих оболочек} // Изв. АН СССР. Сер. матем. -- 1957. -- 


Т.\,21. -- \No\,6. -- С.\,747--784. 


\bibitem{11} 


Велиев М.А., Мамедов Я.Д. {\em Устойчивость решений дифференциальных 


уравнений с нелинейным потенциальным оператором в гильбертовом 


пространстве} // ДАН СССР. -- 1973. -- Т.\,208. -- \No\,1. -- 


С.\,21--24. 


\bibitem{12} 


Лионс Ж.-Л. {\em Некоторые методы решения нелинейных краевых задач}. -- 


М.: Мир, 1972. -- 588 с. 


\bibitem{13} 


Морозов Н.Ф. {\em Исследование колебаний призматического стержня под 


действием поперечной нагрузки} // Изв. вузов. Математика. -- 1965. -- 


\No\,3. -- С.\,121--125. 


\bibitem{14} 


Колмогоров А.Н., Фомин С.В. {\em Элементы теории функций и 


функционального анализа}. Учеб. пособие. -- 5-е изд. -- М.: Наука, 


1981. -- 544 с. 


\bibitem{15} 


Натансон И.П. {\em Теория функций вещественной переменной}. -- М.: 


Наука, 1974. -- 480 с. 


\bibitem{16} 


Лионс Ж.-Л., Мадженес Э. {\em Неоднородные граничные задачи и их 


приложения}. Т.\,1. -- М.: Мир, 1971. -- 372 с. 


\bibitem{17} 


Хартман Ф. {\em Обыкновенные дифференциальные уравнения}. -- М.: Мир, 


1970. -- 720 с. 


 


\end{thebibliography} 


 


\vskip 0.5cm 


 


\post{\it Поволжская академия государственной}{\it Поступила} 


\post{\it службы $($г. Саратов$)$}{27.05.1996} 


 


\label{end} 


\end{document} 


